\documentclass[10.5pt]{article}
\usepackage[margin=2.2cm]{geometry}
\usepackage[T1]{fontenc}
\usepackage{lmodern}
\usepackage{booktabs}
\usepackage{tabularx}
\usepackage{enumitem}
\usepackage{xcolor}
\usepackage[colorlinks=true,linkcolor=blue!50!black,urlcolor=blue!50!black]{hyperref}
\setlist{nosep,leftmargin=1.4em}
\setlength{\parindent}{0pt}
\setlength{\parskip}{4pt}
\emergencystretch=3em

% One entry per change, newest first: \entry{date}{title}
\newcommand{\entry}[2]{\bigskip\noindent{\large\bfseries #2}\hfill{\small\color{black!60}#1}\par\nobreak\vspace{2pt}\hrule\vspace{4pt}}
% Break the page before a field label rather than leave it alone at the bottom.
\newcommand{\keeplines}[1]{\par\vskip 0pt plus #1\baselineskip\penalty-250\vskip 0pt plus -#1\baselineskip}
\newcommand{\field}[1]{\keeplines{3}\textbf{#1}\quad}
\newcommand{\pr}[1]{\href{https://github.com/anurag-r20/Network-Interdiction/pull/#1}{PR~\##1}}
\newcommand{\iss}[1]{\href{https://github.com/anurag-r20/Network-Interdiction/issues/#1}{\##1}}
\newcommand{\f}[1]{\texttt{#1}}

\title{Network Interdiction --- Work Log}
\author{Maintained by Claude Code and Codex for Anurag Ramesh}
\date{Last updated: 26 October 2026}

\begin{document}
\maketitle

\noindent This log records every change to the code, newest first. A new
entry is added, and this PDF rebuilt, each time a piece of code is changed or
implemented. Page~1 is always the current state.

\section*{Current state (25 October 2026)}

\begin{tabularx}{\linewidth}{@{}lX@{}}
\toprule
\textbf{Pull requests} & \\
\pr{20}--\pr{24} & Merged: \pr{20} and \pr{21} on 6 October; \pr{22}, \pr{23} and \pr{24} on 15 October (by bernalde). \\
\pr{25}--\pr{37} & \emph{Open.} All opened by bernalde on 15 October; CI green; no independent review on GitHub. Reviewed locally on 22 October (entry below): none is ready as is. \pr{29}, \pr{33} and \pr{37} need changes; the others need small fixes. Nothing has been posted to GitHub. \\
Merge order & Two stacks: \pr{25}$\to$\pr{27}$\to$\pr{29}$\to$\pr{36}$\to$\pr{37} and \pr{28}$\to$\pr{30}$\to$\pr{31}$\to$\pr{32}$\to$\pr{35}. \pr{26}, \pr{33} and \pr{34} stand alone. \pr{37} already contains \pr{28} and \pr{30}--\pr{35}, so it must not be merged first. \\
\midrule
\textbf{Gurobi} & \emph{Fixed} on 22 October: academic license 2877481 (expires 2027-10-22), bound to the fixed-address adapter \f{eth1}. \\
\midrule
\textbf{Results notebook} & \emph{Rebuilt} on 23 October from merged results only, on the local branch \f{notebook/results-v2} (not pushed). On 25 October, added expected best objective gap against measurement count for the tiny QAOA states; existing figures are unchanged. Sections for open PRs are added as each merges. \\
\midrule
\textbf{Long runs} & None active. \\
Decomposition benchmark & The eight-thread run (14 October, 497 rows, all 27 slide cases) is checked and committed on \f{notebook/results-v2} as \f{results/decomposition\_large.csv}. \\
\midrule
\textbf{Backlog (only on request)} & \\
1 & Follow-up fix for the merged \pr{24}: add the empty-cut probability and a warm-start-only baseline, report $P(\mathrm{opt})$ first, and disclose or drop the $\lambda$ feature. \\
2 & Post the review findings on each open PR (needs your approval). \\
3 & Cap \f{partition\_to\_fit} when no partition fits. \\
4 & LP line, led by the user (decided 23 October): take over \pr{25}'s branch, then study why the LP is exact on these instances, where it breaks, and LP-seeded QAOA. \\
5 & Deferred: implementing arXiv 2610.10772. (Benders, \iss{15}, is now in \pr{27}.) \\
\midrule
\textbf{Needs you} & \\
& Review the rebuilt notebook; then push \f{notebook/results-v2} and open a PR, or ask for changes. \\
& \pr{26}: confirm the MIT license with Hanjing Xu, who wrote most of the original code. \\
& Drop the two old stashes (\f{git stash drop}, twice); I was not allowed to. \\
& Decide whether to commit the IBM hardware results and this log to the repository. \\
\bottomrule
\end{tabularx}

% ---------------------------------------------------------------------------
\section*{Entries}


\entry{26 Oct 2026}{Copy the project to SECQUOIA}
\field{Request}Move the repository to SECQUOIA, or create a new repository
there and copy the project if a transfer is not possible.
\field{Change}GitHub rejected the direct transfer with HTTP 422
(``Repository can't be transferred''). Created the private repository
\href{https://github.com/SECQUOIA/Network-Interdiction}{SECQUOIA/Network-Interdiction}.
The copy includes the current notebook, results, paper, V1 presentation
and worklog on \f{notebook/results-v2}, alongside the source repository's
existing branches and Git history. The current-work branch is the new
default; the source main branch is preserved. Local environments, caches
and ignored build files are excluded.
\field{Scope}Issues and pull requests remain in
\href{https://github.com/anurag-r20/Network-Interdiction}{the original repository};
their historical links in this log still point there. No pull requests
were merged and no experiments were rerun.
\field{Verification}Confirmed that the destination is private and that all
169 project files in the prepared copy match the working checkout byte
for byte. Preserved tracked file modes and rebuilt this PDF.


\entry{26 Oct 2026}{Keep only the V1 presentation}
\field{Request}Remove all other presentation versions, their corresponding
PDFs and scripts, keeping the version ending in \f{\_V1}.
\field{Change}Preserved the author's
\f{Network\_Interdiction\_Project\_Update\_V1.pptx} unchanged and exported
its matching three-page PDF with native PowerPoint. Removed superseded
decks, the old \f{st\_NIP} presentations, PDFs, previews, figure copies,
archived checkpoints, generation metadata and version-specific scripts.
The slide README now points only to V1. Font files and their license remain.
The notebook, source figures, paper and research code are unchanged.
\field{Verification}Checked the V1 package, confirmed the PDF's page count
and titles, and verified that the retained PPTX has the same SHA-256 hash
as before cleanup. V1 is the only remaining project presentation.


\entry{26 Oct 2026}{Use the exact notebook plots in a three-slide presentation}
\field{Request}Replace the separately redrawn slide figures with the same
plots as the notebook. The user selected three slides when asked about
readability and the earlier two-slide limit.
\field{Change}Captured the latest author-edited deck. Slide 1 pairs the
size and sampling figures; slide 2 shows decomposition; slide 3 shows
parameter setting strategies at both $p=1$ and $p=2$. Each embedded PNG
is byte-for-byte identical to its saved notebook figure. No plot is
redrawn, cropped or relabeled. The latest author titles, wording,
formatting and removal of the visible MPS sentence are preserved across
the new layout. The transfer caption and notes clarify which depth its
66\% statistic describes.
\field{Files}The current presentation, PDF and previews use suffix
\f{\_notebook\_plots} in the existing slide directory. The new builder
copies saved figures directly; its manifest records their hashes. Earlier
presentations and all notebook files are untouched.
\field{Verification}Verified all four image payloads against the notebook
PNGs, checked aspect ratios and zero crop, and checked preservation of
author text and unrelated package parts. Native PowerPoint exported three
slides and three PDF pages. All previews were inspected; package,
template and delivery checks pass, with no native text overflow.


\entry{26 Oct 2026}{Add the qubit-size motivation and refresh the presentation plots}
\field{Request}Include the notebook's problem-size figure as a main reason
for decomposition, and update the presentation with the cleaned plot labels.
\field{Change}Started from a new checkpoint of the author's latest saved
deck, including the title ``Decomposition-Based Refinement''. The first
slide now places the full-QUBO and fitted-piece sizes beside the two
decomposition comparisons. It retains all network points, family markers,
the device and simulation limits, and the no-fit warning. Both QAOA plots
are refreshed with concise axis labels, and the SEM bars remain.
The presentation still has two slides. All existing wording, formatting
and notes are preserved; only figure layout and two result-text positions
change. Added figure provenance to the first slide's notes.
\field{Files}The outputs are in \f{slides/network\_interdiction\_update/},
with suffix \f{\_updated}.
Earlier author-edited presentations are untouched. The revision builder,
input hashes, checkpoint and preservation report are stored alongside them.
\field{Verification}Checked that unrelated OOXML parts and existing text
remain unchanged. Rendered both slides in native PowerPoint, inspected the
plots and labels, and ran package, template and PDF-delivery checks.
No notebook data or experiments changed.


\entry{26 Oct 2026}{Concise axis labels throughout the results notebook}
\field{Request}Apply the axis-label cleanup to all notebook plots, without
secondary explanations in the labels.
\field{Change}Removed the QUBO size formula, mean-count parentheticals,
log/symlog notes and decomposition-size ranges from the axes. Kept the
primary quantities, units, and the distinction between the two decomposition
settings. The stacked outcome plots now use ``Probability''.
Refreshed the five affected saved figures and their embedded notebook PNGs.
Data, scales, tick labels, legends, error bars and notebook prose are unchanged.
\field{Verification}Checked that every other notebook field and saved
figure remains unchanged, preserving the user's edits. Ran the existing
notebook checks against the source data.


\entry{26 Oct 2026}{Shorter y-axis label on the sampling-quality plot}
\field{Request}Remove the mean/SEM parenthetical from the y-axis title.
\field{Change}The label is now ``Expected best gap to optimum, \%''.
The SEM bars, data, legend and notebook explanation are unchanged.
Refreshed \f{figures/qaoa\_sample\_quality.png} and only that plot's
embedded PNG in the current notebook, preserving the user's other edits.
\field{Verification}Checked that all other notebook content and saved
figures remain unchanged. The existing sample-quality check verifies the
mean curves and SEM endpoints against the six per-case values.


\entry{26 Oct 2026}{Taller decomposition plot, preserving the author's slide edits}
\field{Request}Explain slide 1's categorical budget-allocation axis and fix
its overly wide, shallow plot while preserving the user's changes.
\field{Change}Captured the latest saved presentation; it was not open in
PowerPoint. Rebuilt only the decomposition image with more vertical space,
kept its aspect ratio, and moved the two introductory lines upward into
space beneath the user's shorter title. The revised PPTX and PDF have
the suffix \f{\_plot\_fixed} in \f{slides/network\_interdiction\_update/}.
The user's original file is untouched. Titles, removed text, formatting,
notes and all of slide 2 are preserved. No notebook or result data changed.
\field{Verification}Package comparison limits changes to slide 1's plot
image and three shape geometries. All other package parts, including
slide 2 and notes, are byte-for-byte unchanged. Checked package integrity
and rendered the revision in native PowerPoint.


\entry{25 Oct 2026}{Two-slide SECQUOIA project update}
\field{Request}Create at most two slides using the SECQUOIA deck, explaining
the supplied project progress report and the notebook's experimental results.
\field{Content}Slide 1 covers the interdiction objective, penalty validation,
partitioning, budget allocation and refinement. It reuses the decomposition
figure to show the trade-off between small pieces and solution quality,
alongside the exact and greedy classical references. Slide 2 pairs the
expected best-of-$N$ gap figure (including SEM) with the depth-two transfer
outcome panel, and summarizes MPS and the separate IBM Fez pilot.
The warm-start claim is qualified: lower QUBO energy did not imply better
interdiction, and much of the probability remained on the empty plan.
Speaker notes explain the cohorts, metric, invalid-plan handling, simulation
limits and hardware provenance.
\field{Files}\f{slides/network\_interdiction\_update/} contains the editable
presentation, matching PDF, slide previews,
builder, native PowerPoint export script, plot copies, and a manifest of
input hashes. The builder uses the neighboring SECQUOIA template v5.19.
Roboto fonts are embedded in the delivered presentation; the font license
is included. The original notebook, figures and experimental data are
unchanged, and no experiments or hardware jobs were run.
\field{Verification}Native PowerPoint exported exactly two slides and two
PDF pages. Package integrity, SECQUOIA publication and delivery checks pass.
Both slide previews were inspected; all 21 native text shapes fit their
boxes within a one-point tolerance. Delivery and layout reports are saved
under \f{qa/}. Both PowerPoint files retain speaker notes, and the PDF
retains figure descriptions; these checks are not an accessibility
certification.


\entry{25 Oct 2026}{SEM bars on the QAOA measurement-count figure}
\field{Request}Add standard-error bars to the new solution-quality figure.
\field{Change}Each point now shows its mean $\pm$ one SEM across the six
cases, using the sample standard deviation ($\mathrm{ddof}=1$) divided by
$\sqrt{6}$. These bars describe variation across cases, not shot noise.
The axis label, notebook explanation and README state this definition.
The saved data and mean curves are unchanged.
\field{Verification}The notebook test independently computes SEM from the
six per-case values and checks both endpoints of every plotted error bar.
All 22 notebook checks pass.


\entry{25 Oct 2026}{QAOA solution quality against measurement count}
\field{Request}Implement the final suggestion in the user's pasted Claude
conversation: add a figure of the best-of-$N$ objective gap against the
number of measurements, with one line per QAOA depth and uniform guessing
as the baseline. Keep the existing outcome figure.

\field{Method}\f{qaoa\_sample\_quality.py} re-simulates the saved tiny-instance
angles, marginalizes auxiliary bits, and scores each interdiction plan by
$f/f^*-1$. An over-budget plan receives the empty plan's score,
$\mathrm{maxflow}/f^*-1$; it cannot earn credit by cutting extra arcs.
An exact survival-tail calculation gives the expected minimum gap after
$N$ independent measurements. Each curve averages those expectations over
the same six cases, with equal weight per case. No angles are retrained,
and no greedy repair is applied. This measures ideal sampling quality,
without training time, execution time, or hardware noise.

\field{Data and notebook}\f{results/qaoa\_sample\_quality.csv} has 288 rows:
six cases, uniform guessing and depths one through three, and twelve
measurement counts from 1 to 4096. Its header records the working-tree
commit, framework versions, and hashes of the inputs and scoring code.
The notebook only reads this CSV. The new figure is
\f{figures/qaoa\_sample\_quality.png}; README explains its metric and
regeneration command. The existing outcome plot and every older figure
are byte-for-byte unchanged. Changes remain local on
\f{notebook/results-v2}; nothing is pushed.

\field{Findings}At ten measurements, mean expected gap is 49\% for uniform
guessing, 19\% at $p=1$, 9\% at $p=2$, and 10\% at $p=3$, reproducing
Claude's scratch calculation. More measurements help; more depth is not
consistently better. With thousands of measurements the gaps approach
zero on these tiny cases, concealing their different sample requirements.
These expectations do not guarantee the result of a particular shot batch.

\field{Verification}All 48 notebook, outcome, and sample-quality tests
pass. Checks include direct enumeration of measurement sequences,
independent NetworkX flow calculations, a Qiskit Statevector cross-check,
coverage and monotonicity of the saved curves, and the plotted means.
Three in-memory defects---rewarding an illegal cut, ignoring $N$, and
reversing the x-bit order---each fail the intended test. The notebook
executes without errors, the new figure was inspected, all eleven input
and source hashes match, and all 21 existing figure hashes are unchanged.


\entry{24 Oct 2026}{QAOA depth axis label spelled out}
\field{Change}At the user's request, the left panel's y-axis now reads
``Probability of optimal plan'' instead of ``P(optimal plan)''; nothing else
changed (commit \f{68b3286}). The notebook executes without errors and
its 21 tests pass.


\entry{23 Oct 2026}{Figure titles removed}
\field{Request}The user asked for no titles on any of the plots.

\field{Change}No figure has a title or suptitle any more. Panel identifiers
moved to axis labels (QAOA depth $p$, decomposition setting) or a caption
under each example network. Each takeaway the titles stated is now a sentence
under its figure, and a figure-size takeaway was added for the problem-size
plot. README and the module docstring say so (commit \f{e4f2f61}).

\field{Verification}The notebook executes without errors and 21 tests pass.
Two are new: one checks that no figure has a title, and one recomputes every
number the text quotes (19 claims) from the raw CSVs; a wrong number planted
in the text made it fail.


\entry{23 Oct 2026}{Gurobi timed at its first optimal incumbent}
\field{Request}The user asked whether the time-against-gap figure used
Gurobi's best-incumbent time from the solver callback, rather than the time
to prove optimality. It did not: it used \f{Gurobi\_time\_sec}, the end of
the solve.

\field{Change}Gurobi is now placed where its MIPSOL callback first reported
the optimal incumbent (\f{Gurobi\_trajectory}). That matches the heuristics,
which are timed until their answer is ready. Across all 27 cases the best
incumbent is the optimum, found at 0.0085--0.048~s; the proof follows within
0.7~ms, so the points barely move. The title, the axis label and the notebook
text say what the times mean (commit \f{68223eb}).

\field{Verification}The notebook executes without errors and its 19 tests
pass. The test recomputes the incumbent times from the trajectory JSON; with
the proof time planted back, it fails.


\entry{23 Oct 2026}{Notebook reloads its figure code}
\field{Problem}The user still saw the unlabelled figure. The notebook on disk
was correct, but Run All in a kernel that had already imported
\f{results\_plots} reuses the old module, so figure edits did not show
without a kernel restart.

\field{Change}The setup cell reloads \f{utils} and \f{results\_plots}, as
the earlier notebook did for \f{utils} (commit \f{f28028d}). The notebook
executes without errors, and its 19 tests pass.


\entry{23 Oct 2026}{Large networks labelled in the problem-size figure}
\field{Request}After the problem-size figure was explained, the user asked
for each large network to be labelled. Three of them (planar 125, ER 90 and
k-NN 100) nearly coincide, and the 50-node ones cluster.

\field{Change}The nine names now sit in a column to the left of the group, in
the same top-to-bottom order as the points, each with a leader line. The
planar 125 label carries the ``no split fits 155 qubits'' note. Commit
\f{9fbdadf} on \f{notebook/results-v2}.

\field{Verification}The notebook executes without errors. 19 notebook tests
pass, including a new check that each network is named once, at its own
point.


\entry{23 Oct 2026}{Notebook names the instances that were run}
\field{Request}The user asked the notebook not to refer to slides or slide
instances, and to say which instances were actually run.

\field{Change}The notebook opens with a table of the four instance sets it
uses: tiny (4 networks, 6 cases), the transfer cohort (6 graphs), mid-size
(9 networks, 23 cases) and large. The large set is the nine networks of
\f{testcases/}: planar with 50, 75 and 125 nodes, Erd\H{o}s--R\'enyi with 50,
70 and 90, and k-NN with 50, 75 and 100, with 288--766 arcs and budgets 4--12
(27 cases). That name replaces ``slide'' in the text, the figure titles and
legends, and the code (\f{load\_large}). \f{results/decomposition\_slide.csv}
became \f{results/decomposition\_large.csv}, with its data bytes unchanged
(commit \f{76df400}).

\field{Verification}The notebook executes without errors and contains no
mention of slides. The 34 notebook, outcome and decomposition tests pass. The
two retitled figures were inspected.


\entry{23 Oct 2026}{Results notebook rebuilt from merged results}
\field{Request}The user found the notebook's plots unhelpful and asked for a
proper results notebook, built first from merged PRs only, with each open PR
added after it merges.

\field{Branch}Local branch \f{notebook/results-v2} from \f{main}, not pushed.
\f{47e1e63} keeps the uncommitted 14 October notebook work unchanged as the
starting point. \f{b012a4e} adds \f{qaoa\_outcomes.py}, \f{6286beb} its
results, \f{adc49c4} the decomposition results, and \f{9478133} the notebook,
\f{results\_plots.py}, tests, README and figures.

\field{New data}\f{results/qaoa\_outcomes.csv} (96 rows) re-simulates every
stored QAOA state exactly and splits it over interdiction plans: optimal,
within budget, empty, or over budget. All 84 stored P(optimal) values
reproduce within $10^{-6}$; it was generated from a clean checkout.
\f{results/decomposition\_slide.csv} is the finished 14 October eight-thread
run (497 rows). Its data bytes are unchanged, and it is checked against the
slide optima, greedy values and max flows.

\field{Notebook}Eight figures, each with a title computed from its data:
example networks with their optimal cuts; QUBO size against simulator and
device limits; a heat map of heuristic gaps; time against gap; QAOA's
outcome split by depth; the parameter-transfer outcome split; within-budget
and optimal MPS shots; and decomposition gaps by budget rule. The notebook
runs in about 15~s and calls no solver. New figure names avoid overwriting
the older figures, which stay in \f{figures/} unused.

\field{Findings}Gurobi proves all 27 slide optima within 0.05~s. Greedy is
optimal on 19 of 27; SA + repair never beats greedy. On the tiny instances,
QAOA's P(optimal) rises from 0.06 to 0.22 by $p=3$, but the empty plan
grows from 0.02 to 0.38, so energy scores mislead. On the held-out cohort,
the transfer methods match target training, and RWS puts 66\% of its
probability on the empty plan at $p=2$. Above 100 qubits, at most 0.4\% of
MPS shots stay within budget. Four pieces with refinement reach the optimum
in 21 of 27 cases; 155-qubit pieces leave a median 22\% gap even with the
best rule (min-cut).

\field{Verification}Notebook tests: 19 pass, checking every figure against
the raw CSVs; four planted bugs in the figure code were each caught. Outcome
tests: 8 pass, and two planted bugs were caught. Decomposition tests: 7
pass. Full suite with the real license: 264 passed, 9 skipped (optional
torch and \f{qaoa\_training\_pipeline}).


\entry{22 Oct 2026}{Independent review of \pr{24}--\pr{37}}
\field{Request}The user asked not to trust the PRs and to review each one
carefully. No repository code was changed.

\field{Method}Nine read-only reviews, each in its own temporary worktree
(all removed afterwards). Each read the PR's diff against its own base,
checked its claims against the linked issue, re-ran its tests and the full
suite with Gurobi's trial license, recomputed committed numbers, and planted
bugs to check that the tests catch them. Nothing was posted or changed on
GitHub. The main checkout was not modified: one stray command in it was
checked and changed nothing.

\field{Verdicts}\par
\begin{tabularx}{\linewidth}{@{}llX@{}}
\toprule
PR & Verdict & Main problem \\
\midrule
\pr{24} (merged) & follow-up fix & Misleading RWS score; transfer no better than one fixed angle set; the $\lambda$ feature uses the exact solution \\
\pr{25} & small fixes & The Lagrangian pass has no time limit \\
\pr{26} & small fixes & Needs the consent of Hanjing Xu, author of most original code \\
\pr{27} & small fixes & Stopping test untested; ignores the thread flag; unconverged runs plotted as converged \\
\pr{28} & small fixes & D-Wave occurrence counts are always 1 after descent \\
\pr{29} & needs changes & The shared price is never applied to the piece solves \\
\pr{30} & small fixes & \f{QCI\_API\_URL} is required, not optional; failures lose finished rows \\
\pr{31} & small fixes & QAOA numbers depend on the machine; one large case aborts the batch \\
\pr{32} & small fixes & No guard against QPU spend; the 300~s timeout loses paid results \\
\pr{33} & needs changes & Large-case embeddings depend on machine speed; 884~KB of regenerable data \\
\pr{34} & small fixes & The Lean proves the arithmetic steps only; docs claim the whole bound \\
\pr{35} & small fixes & XY is no better than a random size-$B$ cut; the table omits that baseline \\
\pr{36} & small fixes & ``Interior'' arcs are not interior; the one hard case is formulation-dependent; data licenses \\
\pr{37} & needs changes & Contains seven other PRs; SA understated; the large QAOA run was never trained \\
\bottomrule
\end{tabularx}

\field{Main finding}The current instances are easy for classical methods.
The LP relaxation equals the optimum on 49 of the 50 slide and mid cases,
and LP rounding plus repair is optimal on all 50 in about 0.02~s each.
\pr{36}'s one ``hard'' case is solved in 0.1~s by a simple exact search
and in about 3.5~s by the standard MILP. The ``hard for Gurobi'' goal of
\iss{1} is not met yet.

\field{Correction}The RWS warm-start score reported for \pr{24} (gap
closed 0.80 and 0.86) mostly reflects probability on the empty cut. That
cut leaves the full flow, the worst answer, yet scores about 0.98. On
$P(\mathrm{opt})$, RWS is worse than standard QAOA at $p=2$ (0.118 against
0.138). GAT and IBM transfer were no better than one fixed angle set.

\field{Tests}Every PR's full suite passes with zero failures. The
remaining skips need optional packages (torch, IBM, QCI) or, at the time,
the full Gurobi license.

\field{Side effects}The \pr{34} review installed the Lean 4.35 toolchain
in \f{\textasciitilde/.elan} and filled the mathlib cache.


\entry{22 Oct 2026}{Gurobi license fixed permanently}
\field{Cause}On WSL2, Gurobi's host ID is the last four bytes of
\f{eth0}'s hardware (MAC) address, and WSL2 gives \f{eth0} a new random
address at every start. Each key therefore stopped working after the next
restart (HostID mismatch on 16 and 20 October).

\field{Fix}Following Gurobi's WSL2 recipe, a boot service
(\f{gurobi-fixed-adapter.service}, set up by the user with \f{sudo}) creates
an unused adapter \f{eth1} with the fixed address \f{02:00:47:52:42:01}.
On the first attempt, udev's default address policy overwrote that
address; the rule \path{/etc/systemd/network/10-gurobi-eth1.link} now
prevents this, and the script waits for udev before setting the address.
The adapter has no IP address and does not affect networking.

\field{License}Installed Gurobi's license tools 13.0.3 in
\f{\textasciitilde/.local/bin}. Activated the new key against \f{eth1}
(host ID \f{47524201}): license 2877481, academic, expires 2027-10-22. The
old file is kept as \f{\textasciitilde/gurobi.lic.bak-20261022}.

\field{Verification}A 3000-variable model, above the trial license's
2000-variable limit, solves. The two tests that need a full license,
\f{test\_lambda\_window} and \f{test\_slide\_results}, now run and pass
(28 passed). Survival across a WSL restart was not tested, at the user's
choice. If the license breaks again, check \f{ip -o link show eth1}; if
the address differs, restart the service. No repository code changed.


\entry{15 Oct 2026}{\pr{22}--\pr{24} merged; thirteen new PRs (recorded 22 Oct)}
\field{Merges}bernalde merged \pr{22}, \pr{23} and \pr{24} into \f{main}
on 15 October (head \f{56e66d6}). \pr{24} had no independent approving
review; it was reviewed afterwards (entry above).

\field{New PRs}bernalde opened \pr{25}--\pr{37} the same day. They cover
LP, Lagrangian and Benders baselines (\iss{15}), the shared budget price
(\iss{14}), real networks (\iss{1}), QCI, Qiskit and IBM backends
(\iss{17}), embedding cost (\iss{4}), penalty bounds with Lean (\iss{2}),
the XY mixer (\iss{3}), raw annealer feasibility (\iss{18}) and the license
(\iss{9}).

\field{Decomposition}The eight-thread decomposition run started on
14 October finished at 21:33 that day: 497 rows covering all 27 slide
cases. It has not been verified or committed.


\entry{14 Oct 2026}{Notebook metrics aligned with quantum preconditioning}
\field{Scope}The user requested a simpler notebook using the same metrics
and plots as the quantum-preconditioning repository. Further experiments
and implementation work require a new user request.

\field{Definitions}Read \f{Gurobi\_QP/notebooks/analysis.ipynb} and its
\f{qp\_gurobi/utils.py}. The notebook now measures the original feasible
interdiction objective, surviving maximum flow $f$, against its exact
optimum $f^*$. It uses $f/f^*-1$ and
$100[1-(f-f^*)/|f^*|]$, matching the reference notebook's executable
definitions. Zero optima are excluded. These are distinct from the earlier
energy scores and normalized flow-reduction plots.

\field{Plots}Reduced the main notebook to three figures: the five-method
objective gap at depths one and two after repair, the classical performance
distribution, and the recorded incumbent trajectory. All saved small
transfer results tie at the optimum after repair. Timing/PAR and scaling
comparisons that need missing trajectories are not reconstructed from final
scores. Existing diagnostic helper regressions remain covered separately.
This change reads saved data and does not run experiments.

\field{Verification}All ten notebook and metric checks pass. They retain
the earlier helper regressions and check the new objective formulas,
missing/zero references, independence from Hamiltonian scores, plotted
trajectory values, and every suboptimal performance point. The notebook
was executed from saved CSVs and its three rendered figures inspected.


\entry{14 Oct 2026}{Depth-one and depth-two comparison for presentation}
\field{Request}Added a side-by-side comparison of all five QAOA methods at
$p=1$ and $p=2$, using the same six held-out graphs and the saved 60-row
statevector experiment. Depths occupy columns; energy gap closed and
optimal-interdiction probability occupy rows. Axis scales match across
depths, graph colors match across panels, and black diamonds mark means.

\field{Artifacts}Updated \f{paper/make\_results.py} to produce a standalone
depth-one figure and a 16:9 comparison in PNG, PDF, and SVG formats under
\f{paper/figures/}, named \f{transfer\_depth\_comparison.*}. The comparison PNG uses
300 dpi. The manuscript retains its depth-two figure. No experiment was
rerun, and no saved measurements were changed.


\entry{14 Oct 2026}{Eight-thread Gurobi policy and full benchmark restart}
\field{Request}The user requested eight threads for every Gurobi run and
then explicitly requested a complete restart for fair timing comparisons.
Stopped the previous decomposition process (PID 108041). Its 153 rows
covering 16 case/setting groups, log, command and checksum are archived in
\path{logs/decomposition-before-eight-threads-20261014-145615/}.
No old timing row will be resumed or combined with the fresh run.

\field{Configuration}A shared model factory explicitly sets
\f{Threads=8} for exact NIP references, penalty checks, Lagrangian solves,
restricted piece models, and the trajectory solver. The benchmark CLI now
defaults to eight instead of one. Its existing explicit override applies
to the dedicated trajectory solve; reference models remain at eight.
The decomposition CSV records its thread budget, refuses to overwrite an
existing output, and logs partition/contraction versus piece-solve phases.
The notebook identifies its historical one-thread timing data separately.

\field{Restart}The fresh run starts at the first of all 27 cases, with both
partition settings and budget refinement. It runs from the clean checkout
\f{/tmp/ni-decomp-eight}, source \f{1ce5438}, as PID 169748, starting at
15:01 EDT. Output, log, and launch metadata are in
\f{logs/decomposition-eight-threads-20261014-150145/}.
Verified \f{Threads=8} on a constructed model before launch; the CSV records
that setting and the source hash. The first two cases produced 20 new rows.
No partial row from the prior run was resumed. The observed initial process
CPU usage exceeded two cores during the piece-solve phase.

\field{CPU interpretation}The previous decomposition models used Gurobi's
automatic thread setting. Observed one-core CPU use was not a confirmed
one-thread Gurobi limit. NetworkX partitioning, exact terminal contraction,
and max-flow calculations remain serial; eight Gurobi threads do not
parallelize those stages.

\field{Verification}The new runtime tests inspect each real model at
\f{optimize()}, including the CLI and nested reference solves. All seven
new checks pass. A scratch mutation forcing one thread fails all six
solver-path checks, while the overwrite-protection test still passes.
All 59 relevant solver, trajectory, decomposition, and penalty-window tests pass. Both \pr{23} CI jobs pass; all \pr{24} checks also pass, including both workflow runs triggered by the base/head updates. Posted and read back the restart summary on \pr{23}. The paper's decomposition protocol now records the eight-thread setting.


\entry{14 Oct 2026}{SIAM manuscript and rendered PDF; transfer results complete}
\field{Paper}Drafted \f{paper/Main.tex} with Anurag Ramesh and David E. Bernal
Neira as the user-confirmed authors. It uses the neighboring quantum
preconditioning project's unchanged SIAM 2017 class and bibliography style.
A separate compatibility file restores standard LaTeX reference hooks for
current hyperref/cleveref. The rendered PDF and Overleaf ZIP are complete; the PDF has eleven pages, with
penalty and contraction proofs, implementation-specific methods, two
result tables, a two-panel figure, limitations, and primary references.
The writing follows the requested approximate ASD-STE100 style while
retaining mathematical terminology. The PDF is a working draft. Verified the rendered title, proofs, tables, figure and references; there are no unresolved references or overflowing text. The Overleaf archive also compiles independently after extraction.

\field{Data}The transfer experiment finished all 60 rows from clean source
\f{ac571aa}: 12 training graphs, six held-out graphs, depths one and two,
five methods. Its largest statevector has 21 qubits ($2^{21}$ complex
amplitudes, 32 MiB at double precision before simulator overhead).
All 28 transfer and saved-artifact checks pass, including recreated
energies, donor choices and a saved GAT on a contracted piece. Four
artifact mutation probes (missing row, wrong energy gap, wrong energy,
and zeroed GAT weights) each fail the intended checks; the untouched
scratch baseline passes all seven artifact tests.

\field{Results}At depth two, mean gap closure is 0.6528 (fixed), 0.6511
(IBM), 0.6621 (GAT), 0.6695 (target-trained), and 0.8610 (warm start).
Warm-start optimal-interdiction probability is lower than target-trained
QAOA (0.1183 versus 0.1377). All repaired outputs are optimal; greedy is
exact because the cohort uses budget one. These results do not demonstrate
an advantage on hard interdiction. The paper uses the common ten-instance
subset for the MPS train/donor comparison, avoiding mismatched cohorts.

\field{GitHub}Committed the raw CSV/JSON/checkpoints, seven artifact tests,
CI command, and result interpretation as \f{d5152ef}; pushed it and the
penalty precision fix \f{da352f5} to \pr{24}. Original CSV line endings are
preserved; the whitespace check passes with CRLF recognized. Checked four
implemented \iss{16} tasks and verified the read-back. The first task
(``land'' QAOA) stays open until \pr{22} merges. All three CI jobs pass at this head. Posted and verified the result summary, then marked \pr{24} ready and read back its state. No PR was merged.

\field{Pending}The original classical decomposition process remains active.
The paper describes its protocol and labels the performance results as
pending; no partial decomposition CSV is used in the manuscript. Worklog,
slides, hardware logs, and the new paper remain local, outside the transfer
PR. The manuscript still needs the full decomposition comparison, harder
instances, multiple seeds, a product-state-only ablation, MPS convergence,
and complete runtime accounting before submission.

\entry{14 Oct 2026}{Transfer PR opened; decomposition result checks prepared}
\field{State}\pr{24} is open on top of \pr{23}. GitHub's core, QAOA and
transfer jobs all passed at \f{ac571aa}. The 12-training/6-test experiment
at depths 1 and 2 is running from \f{/tmp/ni-transfer-benchmark}, with
outputs in \f{/tmp/ni-transfer-output} and log
\f{/tmp/ni-transfer-benchmark.log}. Its provenance records a clean commit.

\field{Decomposition}The original benchmark is still running. Its first
133 rows cover 14 of 54 case/setting combinations (the 50- and 75-node
planar cases). A separate \pr{23} checkout at \f{/tmp/ni-finish-pr23}
holds draft result checks for complete case coverage, ratios, budgets,
qubit limits and monotonic refinement. A copy of the partial table fails
the coverage check, as intended; it is not committed. The active run is
untouched. A new status comment on \pr{23} records the verified handoff.

\field{Result checks}Prepared transfer-table checks for all 60 expected
rows, disjoint graph splits, donor-angle provenance, zero target optimizer
evaluations for transferred schedules, independently enumerated optimal
flows, energy gaps, and reproduction of a held-out energy for each method.
The comparison uses budget one, where greedy is exact; interpretation is
limited to raw quantum energy/probability scores, not a claim of improved
repaired interdictions.

\field{Reference correction}The restricted Gurobi curve now removes MIP
tolerance noise for integer capacities before comparing penalty-window
endpoints, while retaining fractional objectives for fractional-capacity
callers. Added an integer/fractional reference comparison against exhaustive
enumeration; all 21 transfer tests pass. Committed as \f{da352f5}.
Prepared saved-model checks to reproduce the reported donor choice on each
held-out graph. The saved GAT passed a real contracted-piece solve (flow 5 on
\f{tiny\_a}, with only owned arcs cut).

\entry{14 Oct 2026}{Parameter-transfer implementation and verified issue checklists}
\field{What}Checked nine completed tasks across \iss{2}, \iss{4}, \iss{14},
\iss{15}, and \iss{18}, with each updated issue read back from GitHub.
The ``every infeasible assignment is worse than every feasible one'' task
and the closed-form pipeline task remain unchecked: the exact two-sided
$\lambda$ window does not establish those stronger claims.

\field{Implementation}Added flow-network features (capacities, direction,
source/sink, budget, penalty, fixed-arc ownership), the actual IBM
\f{TransferTrainer} adapter, a three-layer GATv2 donor classifier adapted
from the MIS code, and regularized warm starts with a matching tilted mixer.
QAOA now supports restricted interdiction sets for contracted pieces.
Committed as \f{ac571aa}, with optional dependencies and a dedicated
transfer CI job. A clean detached checkout runs the result generation.

\field{Experiment}Added \f{run\_parameter\_transfer.py}: graph-level
training/test separation, fixed-donor, IBM, GAT, per-instance COBYLA and
RWS comparisons; CSV metrics, graph/angle provenance and JSON model
checkpoints. A four-training/two-test, depth-one smoke run completed.
The full comparison will run from a clean committed checkout.

\field{Tests}Before editing code, the current full suite ran with both QAOA
backends required: 181 passed, 17 skipped in 155 seconds. Licensed tests
skipped in the sandbox. New integrations use an isolated temporary Python
environment, preserving the environment of the active benchmark.
The initial 17 transfer tests pass. In a separate scratch copy, six
deliberate defects were each caught: unrestricted repair, wrong mixer
sign, wrong regularizer sign, missing budget feature, farthest-donor
selection, and discarded GAT edge features. The working tree was not
mutated. The full licensed suite passed all 214 tests in 217 seconds;
the final transfer module passed 20 tests in 12 seconds, including the
restricted Gurobi reference and saved-bank reload. The regular environment
without optional transfer dependencies passed 11 tests and skipped five.

\field{Sources}Verified MIS commit \f{ab8434f}; retained its Apache license.
The current IBM commit \f{4bad7bc} uses \f{provide\_params} with depth set
at construction, unlike the older local clone. The dependency is pinned.
Regularized warm starts follow equations (9) and (18) of arXiv:2603.10191,
adapted to this repository's normalized minimization Hamiltonian.

\entry{14 Oct 2026}{Handoff: benchmark fix and mid-size results are pushed}
\field{What}Claude committed \f{0dd3cde}: \f{run\_decomposition\_benchmark.py}
now writes a ``no fit'' row and carries on when a partition cannot meet the
qubit limit; \f{-{}-fit-qubits} sets that limit. Commit \f{bc19e01} adds
\f{results/qaoa\_mid.csv}: 69 donor-angle rows and 30 directly trained rows.
The local branch and GitHub \pr{23} both point to that commit.

\field{Tests}The handoff records two passing benchmark tests and four
passing result-consistency tests, with four deliberate CSV corruptions
caught and the original bytes restored. Codex is verifying current checks
and the remaining results before completing the PR.

\field{State}The restarted decomposition run is active in its existing
detached worktree at \f{0dd3cde}; it is being preserved. The issue-body
edits were prepared but never posted. Only verified completed tasks will
be checked: exact-window computation replaces bisection but still needs
exact solves, so it does not complete the request for instance-only
closed-form bounds. This work log remains local, as before.

% ---------------------------------------------------------------------------
\entry{10 Oct 2026}{Work log started}
\field{What}This document: \f{worklog/worklog.tex}, built into \f{worklog/worklog.pdf} with \f{pdflatex}. Not committed to git.

% ---------------------------------------------------------------------------
\entry{8 Oct 2026}{\pr{23} (draft): mid-size instances, MPS QAOA, better decomposition}
\field{Why}There were no instances between \f{tiny\_c} (22 qubits) and \f{50.planarflow} (912), so nothing beyond the tiny ones could run on a simulator or on a 156-qubit IBM chip.

\field{What}
\begin{itemize}
\item \f{instances.py}: nine seeded networks in \f{testcases/mid/} (planar, Erd\H{o}s--R\'enyi, k-NN), 39--149 qubits. Their exact optima for 23 instance--budget rows are in \f{results/mid\_instances.csv}. Source and sink: the pair with the most arc-disjoint paths.
\item QAOA past 24 qubits: matrix-product-state (MPS) simulation (\f{qaoa\_qiskit.py -{}-method mps -{}-bond-dimension}), a qubit reordering that keeps couplings short (reverse Cuthill--McKee), exact references from Gurobi, scoring from shots, and a new score \f{gap\_closed} (0 = random guessing, 1 = best state). \f{run\_qaoa\_mid.py} runs all 23 rows.
\item Decomposition (\f{decomposition.py}): a \f{marginal} budget split (solve each piece first, give cuts where they save most), \f{iterated} solving, \f{refine\_split()}, \f{partition\_to\_fit()} (fewest pieces that fit a qubit limit), an exact Gurobi solver per piece. \f{run\_decomposition\_benchmark.py} compares every rule with Gurobi's optimum.
\end{itemize}

\field{Tests}193 pass locally; CI green (QAOA job 41 passed; main job 166 passed, 27 skipped). 16 deliberately reintroduced bugs were all caught.

\field{Results so far}
\begin{itemize}
\item \textbf{Mid-size instances are easy classically.} Greedy is optimal on all 23 rows, and on 22 the optimal cut touches the source or sink.
\item \textbf{MPS speed.} The reordering is essential: one 78-qubit evaluation had not finished in 10 minutes without it; with it, 40 qubits take 0.1\,s. At 116 qubits one evaluation takes 15\,s (bond cap 32) or 92\,s (cap 64).
\item \textbf{QAOA at 39--149 qubits} (angles trained on \f{tiny\_a}, reused unchanged):
\begin{center}\small
\begin{tabular}{@{}lccc@{}}
\toprule
depth & energy gap closed (median, range) & shots within budget (median) & optimum after repair \\
\midrule
p = 1 & 0.40 (0.33--0.49) & 0.000 & 18 / 23 rows \\
p = 2 & 0.54 (0.47--0.67) & 0.004 & 18 / 23 rows \\
p = 3 & 0.63 (0.56--0.79) & 0.053 & 18 / 23 rows \\
\bottomrule
\end{tabular}
\end{center}
Training the angles directly on the 30 rows where that was affordable (up to 81 qubits) closes only 0.031 more of the gap (median; range 0.006--0.368). So transferring angles costs little. The catch: almost no raw shot respects the budget (the linear $\lambda$ penalty, \iss{2}); the optimum appears only after greedy repair.
\item \textbf{Decomposition, partial} (each piece solved exactly):
\begin{center}\small
\begin{tabular}{@{}lcccc@{}}
\toprule
instance & pieces & budgets & best variant optimal & marginal, iterated optimal \\
\midrule
\f{50.planarflow} & 4 & 4 & 4 & 2 \\
\f{75.planarflow} & 4 & 3 & 3 & 3 \\
\f{125.planarflow} & 4 & 2 & 2 & 2 \\
\f{50.planarflow} & 8 (fits 156 qubits) & 4 & 3 & 2 \\
\f{75.planarflow} & 30 (fits 156 qubits) & 3 & 2 & 2 \\
\bottomrule
\end{tabular}
\end{center}
Earlier spot checks on k-NN and Erd\H{o}s--R\'enyi instances were mixed: refinement found the optimum on \f{75\_6.2dnnflow} (145) where greedy does not (165), but the \f{marginal} split was much worse than \f{mincut} on \f{70\_0.1.erdosflow} (100 vs 37; optimum 32).
\end{itemize}

\field{Open}Benchmark crash on \f{125.planarflow} (see Current state); result files not yet committed.

% ---------------------------------------------------------------------------
\entry{8 Oct 2026}{IBM Quantum hardware run on \f{ibm\_fez}}
\field{What}Verified access (Open Plan, account \f{open}). Ran the trained QAOA circuits on the real device, 4096 shots each. Cost 8\,s of QPU time (jobs \f{db3pmc2mb58s73886ir0}, 3\,s; \f{db3pmmo4qg6s73c15qc0}, 5\,s), logged in \f{logs/ibm\_jobs.jsonl}. Nothing committed.

\begin{center}\small
\begin{tabular}{@{}lcccc@{}}
\toprule
circuit & 2q gates & P(optimal): hw / ideal & random & approx.\ ratio: hw / ideal \\
\midrule
\f{tiny\_a} B=1, p=1 (14 q) & 41 & 0.257 / 0.260 & 0.125 & 0.789 / 0.807 \\
\f{tiny\_a} B=1, p=2 & 70 & 0.311 / 0.333 & 0.125 & 0.869 / 0.894 \\
\f{tiny\_a} B=1, p=3 & 108 & 0.355 / 0.322 & 0.125 & 0.896 / 0.930 \\
\f{tiny\_d} B=1, p=1 (21 q) & 66 & 0.091 / 0.101 & 0.031 & 0.778 / 0.799 \\
\bottomrule
\end{tabular}
\end{center}
\field{Finding}Finding the optimal cut survives the noise (2--3$\times$ random guessing, close to ideal); the energy degrades a little with depth.

% ---------------------------------------------------------------------------
\entry{7--8 Oct 2026}{CUDA-Q at 21--22 qubits}
\field{What}Re-ran the whole QAOA sweep in CUDA-Q (CPU simulator \f{qpp-cpu}) to check it against Qiskit at the largest tiny sizes, plus a cross-check of CUDA-Q at Qiskit's trained angles.

\field{Results}Cross-check: all 18 rows agree to $5.4\times10^{-7}$. Independent training: p = 1 and p = 2 land within $10^{-4}$ of Qiskit's angles, with identical approximation ratios. p = 3 was cancelled at 3 of 6 cases (all matching). CUDA-Q on CPU is about 80$\times$ slower than Aer (17\,s vs 0.2\,s per evaluation at 21 qubits); its fast simulators need a GPU.

% ---------------------------------------------------------------------------
\entry{7 Oct 2026}{QAOA results and notebook section (\pr{22})}
\field{What}\f{run\_qaoa\_benchmark.py} runs depths 1--3 on the six tiny rows and writes \f{results/qaoa\_tiny.csv}; a new QAOA section in \f{notebooks/analysis.ipynb} with two figures. Also scored the angles trained on \f{tiny\_a}, reused unchanged on the other rows (the transfer baseline).

\field{Results}Approximation ratio 0.80--0.82 at p = 1, 0.92--0.93 at p = 3. P(optimal) relative to random: 6.1--8.5 at p = 3 on the larger B = 1 rows, but only 1.6--2.1 on the B = 2 rows at every depth. Transferred angles lose at most 0.0023 in approximation ratio.

% ---------------------------------------------------------------------------
\entry{6--7 Oct 2026}{\pr{22}: penalty QAOA in Qiskit and CUDA-Q}
\field{What}\f{qaoa\_common.py} (shared Hamiltonian, scoring, training), \f{qaoa\_qiskit.py} (Aer), \f{qaoa\_cudaq.py} (CUDA-Q kernel), \f{requirements-qaoa.txt}, a separate CI job. Both backends are tested against an independent numpy simulation; 17 reintroduced bugs were all caught.

\field{Results}On \f{tiny\_a} B=1, p = 2: P(optimal) 0.333 vs 0.125 for random guessing; both backends train to the same angles. At B = 2, lower energy does not mean better cuts.

% ---------------------------------------------------------------------------
\entry{6 Oct 2026}{Research: GAT and TransferTrainer (no code)}
\field{Findings}IBM's \f{TransferTrainer} is a lookup (features $\to$ stored angles), not a learned model; it can serve NIP with an NIP feature set, a nearest-match rule and no angle averaging. The MIS GAT work classifies a graph into a cluster of donor graphs, then tries the nearest donors' angles; it handles only unweighted graphs at p = 2. Neither is implemented yet.

% ---------------------------------------------------------------------------
\entry{6 Oct 2026}{P0: merges, slide fix, proposal sentence}
\field{What}Merged \pr{20} (\f{c08a5f7}) and \pr{21} (\f{48c8581}, after resolving one conflict). Corrected slide 2 of the deck: max flow 22 (not 25), reduction 77\% (not 80\%), saved as \f{st\_NIP\_corrected.pptx}. Posted the corrected proposal results sentence on \iss{18}.

% ---------------------------------------------------------------------------
\entry{6 Oct 2026}{Repository review and priority list}
\field{What}Read the code, issues and PRs; ranked next steps (P0 unblock, P1 hard instances and baselines, P2 experiments, P3 blocked items).

\end{document}
